\scanpage{133}
et \deleted{\(\unsure{F}\)} par lemme 1 il suffit de regarder les \uline{extrêmes}.
Celui de droite est O.K. \deleted{\unsure{connu}} par \uline{définition},
celui de gauche par hyp. de récurrence. O.K.

\uline{Proposition 4}\quad
Soit \(i:U\mathbin{\text{---}}X\) une immersion ouverte,
\inserted{sur \(X\) excellent, de dim \(\le d\),}
\(F\) un \deleted{\unsure{le}} \(Q_l\)-faisceau
\deleted{\(\unsure{d}\)}\inserted{\(d'\)}-admissible.
Alors les \(R^ni_*(F)\) sont
\(\unsure{(d+d'\deleted{2}\inserted{1})|(d+d'-1)}\)-admissibles.
\figure{\begin{tikzpicture}\draw (0,2)--(0,0);\end{tikzpicture}}

Récurrence sur \(d+d'\),
\deleted{En effet, soit \inserted{\unsure{de façon que}}
\(\unsure{W|V}\subset U\), \(\unsure{W|V}\) dans \(U\), tel que
\(\unsure{F\vert W|F\vert V}\) \illegible{1}.}
\inserted{trivial si \(d+d'=0\) donc \(d=0,d'=0\).
Supposons \(d+d'>0\), et \deleted{\unsure{est}} \unsure{prouvé} si
\(\delta+\delta'<d+d'\).}

On peut supposer \(X\) affine donc \unsure{noeth.|exclt},
et par réc. noeth., \deleted{supposons} le th. est démontré pour les
\(Y\subset X\) fermés, \(Y\ne X\), et \(X\ne\varnothing\).

Soit donc \(V\underset{j}{\subset}U\) ouvert dense de \(U\), régulier, alors
\diagram{\[
\begin{tikzcd}[row sep=small,column sep=small]
&Q^*\arrow[dl]\\
F\arrow[rr,no head]&&\unsure{\mathbb R|\mathbf R}j_*F\arrow[ul]
\end{tikzcd}
\]}
où \(Q\) est concentré sur \(Z\cap U\),
\(Z=\overline{U-V}\) \uline{\unsure{vu|pris}} dans \(X\), donc \(\ne X\).
Donc on est ramené à prouver que les
\(R^ni_*(Rj_*(F))\) et \(R^ni_*(Q^*)\) sont admissibles.
Pour le
% archive foliation: 132
